\documentclass[10pt,letterpaper]{article}
\usepackage{url}
\usepackage[colorlinks,urlcolor=blue]{hyperref}
\usepackage[UTF8]{inputenc}
%\usepackage[latin1]{inputenc}
\usepackage[T1]{fontenc}
\usepackage{uarial}
\renewcommand\familydefault{\sfdefault}
\usepackage{amsmath,amssymb,amsfonts,latexsym,cancel, wrapfig, multicol,enumerate,graphicx, xcolor,paralist,tikz,tkz-tab,dsfont}
\usepackage[letterpaper,showframe=false]{geometry}
\usepackage[spanish]{babel}
\usepackage{color}
\usepackage{hyperref}
%colores------------------------------------------------
\definecolor{encabezado}{RGB}{64,107,204} % % % Principal
\definecolor{titulos}{RGB}{112,90,45}
\definecolor{actividades}{RGB}{112,90,45}
\definecolor{cuadro1}{RGB}{204,219,226}
\definecolor{cuadro2}{RGB}{226,210,178} % % % tercero
\definecolor{cuadro3}{RGB}{255,230,179}
\definecolor{CRecapitulacion}{RGB}{204,219,226}
\definecolor{CcontenidosObjetivos}{RGB}{204,219,226} % % Secundario
\definecolor{ContenidosObjetivos}{RGB}{112,90,45}
\definecolor{numerosTabla}{RGB}{0,0,0}


%--------------- Colores IS22-----------------------
%\definecolor{encabezado}{RGB}{0,134,75} % % % Principal
%\definecolor{titulos}{RGB}{1,110,111}
%\definecolor{actividades}{RGB}{1,110,111}
%\definecolor{cuadro1}{RGB}{173,211,144}
%\definecolor{cuadro2}{RGB}{249,241,133} % % % tercero
%\definecolor{cuadro3}{RGB}{253,199,117}
%\definecolor{CRecapitulacion}{RGB}{188,228,230}
%\definecolor{CcontenidosObjetivos}{RGB}{188,228,230} % % Secundario
%\definecolor{ContenidosObjetivos}{RGB}{1,110,111}
%\definecolor{numerosTabla}{RGB}{70,70,70}

%--------------- Colores antes de IS22-----------------------
%\definecolor{encabezado}{RGB}{65,130,185}
%\definecolor{titulos}{RGB}{65,130,185}
%\definecolor{actividades}{RGB}{239,80,0}
%\definecolor{cuadro1}{RGB}{216,216,216}
%\definecolor{cuadro2}{RGB}{249,206,17}
%\definecolor{cuadro3}{RGB}{239,159,77}
%\definecolor{CRecapitulacion}{RGB}{216,216,216}
%\definecolor{CcontenidosObjetivos}{RGB}{138,186,212}
%\definecolor{ContenidosObjetivos}{RGB}{86,86,86}
%\definecolor{numerosTabla}{RGB}{48,103,133}
%-------------------------------------------------------

\DeclareGraphicsExtensions {.pdf,.png,.jpg}

%Dimensiones-------------------------------------------
\textheight=26cm
\textwidth=20 cm
\topmargin=-3.2 cm
\oddsidemargin=-2 cm
\evensidemargin=0 cm
\parskip=0 mm
\parindent=0 mm
%-----------------------------------------------------


%\usepackage[T1]{fontenc}
\usepackage{uarial}
\renewcommand\familydefault{\sfdefault}
%\usepackage{blindtext}

\newcommand{\URLLibroMN}{\href{https://ezproxy.itcr.ac.cr/login?url=https://www.bibliotecaebook.com/my-library}{\textbf{Acceso al e-book}} Libro de Métodos Numéricos. Autor Jeffry Chavarría Molina.}

\newcommand{\CursoOctave}{\href{https://lugfi.github.io/curso-octave/index.html}{\textbf{Ingreso al curso}}}

\newcommand{\PaqueteSymbolic}{\href{https://docs.sympy.org/latest/tutorials/intro-tutorial/index.html}{\textbf{Paquete Sympy:}} Paquete para realizar cálculo simbólico en Python}
%--------------------------------
%Videos


\newcommand{\URLVideoUno}{\href{https://youtu.be/_7Aw-ay67DM}{\textbf{Vídeo 1:}} Integración con datos no igualmente espaciado}

\newcommand{\URLVideoDos}{\href{https://tecnube1-my.sharepoint.com/:v:/g/personal/ccalderon_itcr_ac_cr/EfR-z83flapBtUcMU36DMr4BlEX-YLvYJBM7eMtV8kM4Mw?e=9Ed7WD}{\textbf{Video 2:}} Cuadratura Gaussiana. Introducción}

\newcommand{\URLVideoTres}{\href{https://tecnube1-my.sharepoint.com/:v:/g/personal/ccalderon_itcr_ac_cr/EZJCmoE_ZydJqdjObEAX0EcBSlz-IyN2drdlT4QuWE759g?e=3r3SLz}{\textbf{Video 3:}} Cuadratura Gaussiana. Ejemplo}

\newcommand{\URLVideoCuatro}{\href{https://tecnube1-my.sharepoint.com/:v:/g/personal/ccalderon_itcr_ac_cr/ETdNUpQs4Q5LjHNWW62PxP8BUNPz02wIHSAs_nfhxJRW2g?e=5s6yyh}{\textbf{Video 4:}} Gauss-Legendre}

\newcommand{\URLVideoCinco}{\href{https://youtu.be/j0SyekJzuvg}{\textbf{Video 5:}} Cuadratura Gaussinaa en Octave}


%------------Vídeos teoricos---------------
\newcommand{\URLVideoTeoriaActiI}{\href{https://drive.google.com/file/d/1Is-lhFgiC0lT138Ikfc-pgqc2trrvBBq/view?usp=sharing}{\textbf{Vídeo de teoría 1:}}  Complementario al libro}
\newcommand{\URLVideoTeoriaActiII}{\href{https://drive.google.com/file/d/1kloExvEhsScDhAK6WMXRjhp5rQ62zYDm/view?usp=sharing}{\textbf{Vídeo de teoría 2:}} Complementario al libro}
\newcommand{\URLVideoTeoriaActiIII}{\href{https://drive.google.com/file/d/1H7o1l6RkJaNmG_7-_cLjfnSolojOus38/view?usp=sharing}{\textbf{Vídeo de teoría 3:}} Complementario al libro}


%---------Infografías-----------------



%----------Solucionario-------------
\newcommand{\FolletoPractica}{\href{https://tecnube1-my.sharepoint.com/:b:/g/personal/jchavarria_itcr_ac_cr/ESbC_k0vPWdMovfv0-d9GpoBcM31nPWa26ScY419frpBdw?e=hutomi}{\textbf{Folleto de práctica}} Folleto completo de práctica del curso.}

\newcommand{\FolletoPracticaDOS}{\href{https://tecnube1-my.sharepoint.com/:b:/g/personal/jchavarria_itcr_ac_cr/ESbC_k0vPWdMovfv0-d9GpoBcM31nPWa26ScY419frpBdw?e=hutomi}{\textbf{Folleto de práctica}}}


%----------Podcast---------------------------------------------

\newcommand{\URLpodcast}{\href{https://soundcloud.com/angie-solis-palma/logica-matematica-y-conectivas-logicas?utm_source=clipboard\&utm_campaign=wtshare\&utm_medium=widget\&utm_content=https\%253A\%252F\%252Fsoundcloud.com\%252Fangie-252Fsol-palmamama}{\textbf{Podcast:}} Lógica matemática.}



\newcommand{\URLFacebook}{\href{https://tecdigital.tec.ac.cr/dotlrn/catedras.tec/matematicadiscreta/one-community?page_num=0&loc=}{enlace}}

\newcommand{\URLPortalComunidad}{\href{https://tecdigital.tec.ac.cr/dotlrn/catedras.tec/matematicadiscreta/one-community?page_num=0&loc=}{portal de la comunidad del TecDigital para MD}}

\newcommand{\URLPortalGrupoTecDigital}{\href{https://tecdigital.tec.ac.cr/dotlrn/classes/MA/MA1403/S-2-2022.CA.MA1403.6/one-community?page_num=0&loc=}{portal del TecDigital para el grupo 06}}

\newcommand{\URLConsulta}{\href{https://tecdigital.tec.ac.cr/dotlrn/catedras.tec/matematicadiscreta/one-community?page_num=5&loc=}{enlace}}

\newcommand{\ImagenDM}{\includegraphics[scale=0.7]{imagenes/MN.png}}

\newcommand{\Motivacion}{\includegraphics[scale=0.6]{imagenes/motivacion.png}}


\begin{document}
\setlength{\parskip}{0.3cm}

%Encabezado-----------------------------------------------------------------------------------	






\begin{center}
{ \fboxsep 12pt
	\fcolorbox {encabezado}{encabezado}{
	\begin{minipage}[t]{19.5cm}
		{\color{white} 
		TEC | MÉTODOS NUMÉRICOS GRUPO SEMIPRESENCIAL \hfill  II SEMESTRE 2026
 \\
		\rule[0.1cm]{19.5cm}{0.01cm}
		\begin{center}
			{\Huge  \textbf{Semana 11}}
		\end{center}
		\vspace{-0.5cm}
		\rule[0cm]{19.5cm}{0.01cm}
		\begin{center}
			{  \textbf{Integración Numérica}}
		\end{center}
					}
	\vspace{-0.5cm}
	\end{minipage}
	} 
}
\end{center}


%Pagina1-------------------------------------------------------------------------
\begin{minipage}[c]{14cm}
	\begin{center}
		\includegraphics[height=5.4cm,width=11cm]{imagenes/estudiantes}
	\end{center}
		
{\color{titulos} {\large \textbf{Descripción para la semana:}}}\\
		
Durante la \textbf{semana 11} se repasó la integral definida, su interpretación gráfica, el teorema fundamental del cálculo y la regla de Barrow. Se introdujeron las fórmulas de Newton-Cotes y se estudiaron la \textbf{regla del rectángulo simple y la regla del rectángulo compuesta}. Esta orientación organiza el repaso y la práctica de esos contenidos.

{\color{titulos} {\large \textbf{ Modo de trabajo:}}}
		
{ \fboxsep 12pt
	\fcolorbox {cuadro1}{cuadro1}{
		\begin{minipage}[l][9.6cm]{3.15cm}
			\begin{center}
				\includegraphics[scale=0.2]{imagenes/uno}
			\end{center}
			{\color{ContenidosObjetivos} {\large Material a utilizar}} \\ \\
			{\footnotesize
\href{https://jchavarriam-pixel.github.io/ITCR_CURSOS/MN/IntegracionNumerica/}{\textbf{Guía web de integración numérica}}\\[0.2cm]
Videos, explicaciones escritas, fórmulas, cálculos y resultados.\\[0.2cm]
Consulten también las notas de clase enlazadas en la página.
} 
\vspace{0.5cm}
						
		\end{minipage}
		}
}
{ \fboxsep 12pt
		\fcolorbox {cuadro2}{cuadro2}{
			\begin{minipage}[l][9.6cm]{3.9cm}
				\begin{center}
					\includegraphics[scale=0.2]{imagenes/dos}
				\end{center}
				{\color{ContenidosObjetivos} {\large Vídeos}} \\ \\
				{\footnotesize
\textbf{Recursos de estudio:}\\[0.2cm]
Notas de clase y libro de texto para el repaso de integración definida y la introducción a Newton-Cotes.\\[0.2cm]
\href{https://jchavarriam-pixel.github.io/ITCR_CURSOS/MN/IntegracionNumerica/\#rectangulo-compuesto}{\textbf{Video 01 de la guía web:}} Regla compuesta del rectángulo.\\[0.2cm]
Intenten los cálculos antes de consultar el desarrollo escrito.
}
\end{minipage}
	}
}
{ \fboxsep 12pt
	\fcolorbox {cuadro3}{cuadro3}{
		\begin{minipage}[l][9.6cm]{3.15cm}
			\begin{center}
				\includegraphics[scale=0.2]{imagenes/tres}
			\end{center}
			{\color{ContenidosObjetivos} {\large Evaluación}} \\ \\
			 \footnotesize 
			 \textbf{Práctica de la semana:}

Resuelvan los ejercicios de integral definida y de las reglas del rectángulo simple y compuesta indicados en esta orientación.

Conserven sus procedimientos y registren las dudas para consultarlas.
\vfill
\end{minipage}
	} 
} 
\end{minipage}\hspace{0.2cm}\begin{minipage}[c]{5cm}
\begin{center}
{ \fboxsep 12pt
		\fcolorbox {CcontenidosObjetivos}{CcontenidosObjetivos}{
\begin{minipage}[l][21cm]{5cm}
\begin{center}
	\ImagenDM
\end{center}
					
{\color{ContenidosObjetivos}\textbf{ \small Contenidos:}}\\
					
{\footnotesize
{ \tiny $\blacksquare$} Métodos numéricos para la integración de funciones \\ \\
{ \tiny $\blacksquare$}  Programación de métodos para la integración de funciones\\
}

{\color{white}\rule[0.1cm]{5cm}{0.01cm}}
		{\color{ContenidosObjetivos} \textbf{ \small Objetivos:}}\\
					
{\footnotesize
{ \tiny $\blacksquare$} Estudiar métodos asociados la integración de funciones \\ \\
{ \tiny $\blacksquare$} Programar los métodos para
la integración de funciones
}

{\color{white}\rule[0.1cm]{5cm}{0.01cm}}\\

%\vspace{2cm}
%\hspace{-7pt}\begin{minipage}[4cm]{5cm}
%\begin{center}
%\includegraphics[scale=0.14]{imagenes/QRFacebook}
%
%Comunidad de la cátedra de Matemática Discreta en el tecDigital (\URLFacebook).
%\end{center}
%\end{minipage}

\end{minipage}
}}
\end{center}
\end{minipage}	 
%----------------------------------------------------------------------------


\newpage

%Página2----------------------------------------------------------------------
\begin{multicols}{2}
{\color{actividades}\textbf{Actividad 1: Integral definida}\par}
\begin{itemize}
\item Repasen la definición de integral definida y la partición de un intervalo. Consulten la sección 5.3 del libro, páginas 339--342.
\item Interpreten la integral como área con signo y distingan este valor del área geométrica total.
\item Repasen el teorema fundamental del cálculo y la regla de Barrow: si $F'=f$, entonces $\int_a^b f(x)\,dx=F(b)-F(a)$, bajo las hipótesis correspondientes.
\end{itemize}

{\color{actividades}\textbf{Actividad 2: Newton-Cotes y rectángulo}}
\begin{itemize}
\item Revisen la idea de las fórmulas de Newton-Cotes: aproximar la función mediante un polinomio interpolador e integrar ese polinomio.
\item Estudien la regla del rectángulo simple y su interpretación gráfica, según la convención utilizada en clase.
\item Estudien los ejemplos 123, 124 y 125 del libro relacionados con la regla del rectángulo.
\end{itemize}

\columnbreak
{\color{actividades}\textbf{Actividad 3: Rectángulo compuesto}\par}
\begin{itemize}
\item Construyan una partición uniforme con $n$ subintervalos y ancho $h=(b-a)/n$.
\item Identifiquen los puntos de evaluación utilizados en clase y sumen las contribuciones de los rectángulos.
\item Revisen el video 01 de la guía web y su desarrollo escrito. En este recurso se emplean extremos izquierdos.
\item Comparen el resultado con una integral de referencia y distingan el error real de una cota del error.
\end{itemize}

{\color{actividades}\textbf{Actividad 4: Ejercicios de práctica}}

Los ejercicios 1--3 se elaboraron para esta orientación; el ejercicio 4 corresponde al video 01 de la guía web. Las respuestas se incluyen en la página siguiente.
\begin{enumerate}
\item Para $f(x)=x^2$ en $[0,2]$, dibujen la región y calculen la integral definida mediante la regla de Barrow. Expliquen su relación con el área.
\item Para $f(x)=x-1$ en $[0,2]$, calculen la integral definida y el área geométrica total. Justifiquen por qué difieren.
\item Para $\int_0^2 x^2\,dx$, apliquen la regla del rectángulo simple con extremo izquierdo. Luego, apliquen la regla compuesta con extremos izquierdos para $n=4$ y $n=8$. Comparen los errores absolutos con el valor exacto.
\item Resuelvan el enunciado del video 01: aproximen $\int_1^2 \frac{\operatorname{sen}x}{x}\,dx$ con rectángulos izquierdos y determinen $n$ para garantizar un error absoluto menor que $10^{-5}$. Comparen después con la solución de la guía.
\end{enumerate}
\end{multicols}

{ \fboxsep 12pt
\fcolorbox {CRecapitulacion}{CRecapitulacion}{
\begin{minipage}[c]{12cm}
{\color{ContenidosObjetivos}\large {\large Recapitulación}:}\\
		
Se les recuerda la importancia de completar las actividades asignadas de la sección virtual para un desarrollo adecuado de su aprendizaje. Completen las actividades de repaso de la semana 11 y conserven sus cálculos para repasarlos y consultar sus dudas. \\
		
Se le recuerda que no está solo en el proceso. Puede realizar sus dudas en el grupo de Teams.\\
		

\end{minipage}
		\begin{minipage}[c]{7cm}
			\begin{center}
				\Motivacion
			\end{center}
			
		\end{minipage}
	}
}


\newpage
{\color{titulos}\Large\textbf{Respuestas de la actividad 4}}

Intenten resolver cada ejercicio antes de consultar estas respuestas. En las reglas del rectángulo se utilizan extremos izquierdos, tal como indican los enunciados.

{\color{actividades}\textbf{1. Integral de $x^2$ en $[0,2]$}}

Una primitiva es $F(x)=x^3/3$. Por la regla de Barrow,
\[
\int_0^2 x^2\,dx=\left[\frac{x^3}{3}\right]_0^2=\boxed{\frac83}.
\]
La región está entre la curva $y=x^2$, el eje horizontal y las rectas $x=0$ y $x=2$. Como $x^2\geq0$, la integral coincide con el área geométrica: $8/3$ unidades cuadradas.

{\color{actividades}\textbf{2. Integral de $x-1$ y área geométrica}}
\[
\int_0^2(x-1)\,dx=\left[\frac{x^2}{2}-x\right]_0^2=\boxed{0}.
\]
La función cambia de signo en $x=1$. Las dos regiones son triángulos de base y altura iguales a 1: cada uno tiene área $1/2$. La integral suma las áreas con signo, $-1/2+1/2=0$, mientras que el área geométrica total es $\boxed{1}$ unidad cuadrada.

{\color{actividades}\textbf{3. Rectángulos para $\int_0^2 x^2\,dx$}}

Con un solo rectángulo izquierdo, $R=(2-0)f(0)=\boxed{0}$. Para la regla compuesta,
\[
L_n=h\sum_{i=0}^{n-1}(ih)^2,\qquad h=\frac2n,
\qquad L_n=\frac{4(n-1)(2n-1)}{3n^2}.
\]
\begin{center}
\begin{tabular}{c|c|c|c}
Método & $h$ & Aproximación & Error absoluto $|8/3-L_n|$\\ \hline
Simple izquierdo & $2$ & $0$ & $8/3\approx2.666667$\\
Compuesto, $n=4$ & $1/2$ & $7/4=1.75$ & $11/12\approx0.916667$\\
Compuesto, $n=8$ & $1/4$ & $35/16=2.1875$ & $23/48\approx0.479167$
\end{tabular}
\end{center}
Las aproximaciones son menores que la integral porque $x^2$ es creciente y se usan extremos izquierdos. Al aumentar el número de rectángulos, el error disminuye.

{\color{actividades}\textbf{4. Rectángulo compuesto del video 01}}

Para $f(x)=\operatorname{sen}(x)/x$ en $[1,2]$, con ángulos en radianes, el desarrollo de la guía establece $|f'(x)|\leq M=0.4354$. La cota para rectángulos izquierdos da
\[
|I-L_n|\leq\frac{(b-a)^2M}{2n}=\frac{0.4354}{2n}.
\]
Para garantizar un error menor que $10^{-5}$ mediante esta cota, se exige $n>21770$; por tanto, el menor entero que satisface esta condición es $\boxed{n=21771}$. Entonces,
\[
L_{21771}=\frac1{21771}\sum_{i=0}^{21770}
\frac{\operatorname{sen}(1+i/21771)}{1+i/21771}
\approx\boxed{0.659338790}.
\]
La cota es aproximadamente $9.99954\times10^{-6}<10^{-5}$. Con $I\approx0.659329906$, el error real es aproximadamente $8.88\times10^{-6}$, y respeta la cota.



%{\color{actividades} \textbf{   $\bullet$ Códigos QR de los vídeos:}}\\
%\begin{center}
%\begin{tabular}{|c|c|c|c|}
%\hline 
%\textbf{¿Para que sirven las matemáticas?} & \textbf{Motivación a la lógica} & \textbf{Vídeo de teoría 1} &\textbf{Vídeo de teoría 2} \\ \hline &&&\\  
%\includegraphics[scale=0.1]{imagenes/QRMoti1} & \includegraphics[scale=0.1]{imagenes/QRMoti2} & \includegraphics[scale=0.1]{imagenes/QRVideo1} & \includegraphics[scale=0.1]{imagenes/QRVideo2}\\ 
%\hline 
%\textbf{Vídeo 3}&\textbf{Vídeo 4}&\textbf{Vídeo 5}&\\ 
%\hline &&&\\  
% \includegraphics[scale=0.1]{imagenes/QRVideo3} & \includegraphics[scale=0.1]{imagenes/QRVideo4} & \includegraphics[scale=0.1]{imagenes/QRVideo5} & \\ \hline
%\end{tabular} 
%\end{center}
\end{document}

